% Chapter 5: streaming

\chapter{Streaming}

\label{Chapter5}

\file{fluxlb/core/streaming.py} holds one function, \code{stream()}: the push of every
post-collision population one lattice spacing along its own velocity, with periodic wrap.
It is the half of the update that is not collision, and the reason the collision seam of
\file{CLAUDE.md} can be drawn where it is. \Cref{sec:stream-theory} gives the theory and
\cref{sec:stream-impl} the implementation notes.

%----------------------------------------------------------------------------------------

\section{Theory}\label{sec:stream-theory}

Streaming is a pure shift: every post-collision population $f_i^{*}(\mathbf{x}, t)$ is moved
one lattice step along its own velocity,
\begin{equation}
f_i(\mathbf{x} + \vc_i\Delta t,\; t + \Delta t) = f_i^{*}(\mathbf{x}, t).
\label{eq:stream}
\end{equation}
It is exact, linear, non-local and touches no values; it only relocates them
\parencite{Kruger2017}. In PyTorch that is a \code{torch.roll} per direction (periodic by
construction), with boundaries then overriding the populations that crossed a wall
(\cref{fig:push}). Because it is a permutation of the population array, it is also the
operation that becomes a permutation unitary on the quantum track.

\begin{figure}[h]
\centering
\begin{adjustbox}{max width=\linewidth}% Push streaming on D2Q9 and the periodic wrap of torch.roll.
\begin{tikzpicture}[
  >={Stealth[length=2.2mm]},
  node/.style={circle, fill=gray!25, draw=gray!60, inner sep=0pt, minimum size=6pt},
  pop/.style={->, line width=0.9pt, blue!60!black},
  hi/.style={->, line width=1.5pt, red!70!black},
  lbl/.style={font=\small, inner sep=1pt},
  note/.style={font=\small, gray!80!black, align=left},
  cell/.style={draw, fill=blue!7, minimum size=9mm, font=\small\ttfamily},
]
% ---- left: push from one node to its eight neighbours ----
\begin{scope}[shift={(0,0)}]
  \draw[gray!30, step=2.4] (-2.5,-2.5) grid (2.5,2.5);
  \foreach \x in {-2.4,0,2.4} \foreach \y in {-2.4,0,2.4} \node[node] at (\x,\y) {};
  \node[node, fill=black!70] (x) at (0,0) {};
  \node[lbl, below left=2pt] at (x) {$\mathbf{x}$};
  % seven ordinary populations leave the node; label sits beside the arrow tip
  \foreach \i/\dx/\dy in {2/0/1, 3/-1/0, 4/0/-1, 5/1/1, 6/-1/1, 7/-1/-1, 8/1/-1} {
    \draw[pop] (0,0) -- ({2.4*\dx - 0.3*\dx}, {2.4*\dy - 0.3*\dy});
    \node[lbl, blue!60!black] at ({1.55*\dx + 0.42*\dy}, {1.55*\dy - 0.42*\dx}) {$f_{\i}^{*}$};
  }
  % the highlighted one: f_1^*(x) becomes f_1(x + c_1)
  \draw[hi] (0,0) -- (2.1,0);
  \node[lbl, red!70!black, above] at (1.1,0.08) {$f_1^{*}(\mathbf{x},t)$};
  \node[node, fill=red!70!black] at (2.4,0) {};
  \node[lbl, red!70!black, align=center, anchor=north] at (2.4,-0.25)
    {$f_1(\mathbf{x}+\vc_1\Delta t,\, t+\Delta t)$};
  % rest population stays
  \node[lbl, black!70, anchor=north east] at (-0.15,-0.35) {$f_0^{*}$ stays};
  \node[font=\small\bfseries, align=center] at (0,-4.0) {push: every $f_i^{*}$ leaves along its own $\vc_i$};
\end{scope}

% ---- right: one direction on one axis, periodic wrap ----
\begin{scope}[shift={(6.3,1.0)}]
  \node[note, anchor=south west, font=\small\ttfamily] at (-0.4,1.05) {torch.roll(f[1], shifts=(+1, 0), dims=(0, 1))};
  % before
  \foreach \k in {0,...,4} {
    \node[cell] (a\k) at (\k*1.0,0.4) {\k};
  }
  \node[lbl, left=4pt of a0] {$f_1^{*}$};
  \node[lbl, right=4pt of a4] {$x \rightarrow$};
  % after
  \foreach \k in {0,...,4} {
    \pgfmathtruncatemacro{\src}{mod(\k+4,5)}
    \node[cell] (b\k) at (\k*1.0,-1.6) {\src};
  }
  \node[lbl, left=4pt of b0] {$f_1$};
  \foreach \k in {0,...,3} {
    \pgfmathtruncatemacro{\kk}{\k+1}
    \draw[pop] (a\k.south) -- (b\kk.north);
  }
  % wrap: last cell goes to first
  \draw[hi, rounded corners=3pt] (a4.south) -- ++(0,-0.45) -- ++(-4.0,0) -- (b0.north);
  \node[lbl, red!70!black, anchor=west] at (4.55,-0.55) {wrap};
  \node[note, anchor=north west, text width=6.4cm] at (-0.4,-2.35)
    {Cell $x$ holds the value that came from $x-1$; the value leaving the last cell re-enters
     at the first. A wall boundary overwrites that entry afterwards.};
  \node[font=\small\bfseries] at (2.0,-5.0) {periodic by construction};
\end{scope}
\end{tikzpicture}
\end{adjustbox}
\caption{Push streaming. Left: on D2Q9 every post-collision population leaves its node along
its own velocity; $f_1^{*}$ at $\mathbf{x}$ becomes $f_1$ at $\mathbf{x} + \vc_1\Delta t$,
and the rest population stays. Right: one direction on one axis as \code{torch.roll} sees
it. The value leaving the last cell re-enters at the first, which is what makes the step
periodic by construction.}
\label{fig:push}
\end{figure}

$f^{\mathrm{eq}}$ never appears in streaming. It lives entirely in collision, which is
local: at each node you take moments of the current populations, \cref{eq:moments}, build
$f^{\mathrm{eq}}(\rho, \vu)$ from them, \cref{eq:hermite-terms}, and relax $f_i$ toward it,
\begin{equation}
f_i^{*} = f_i - \frac{1}{\tau}\left(f_i - f_i^{\mathrm{eq}}\right).
\label{eq:bgk-relax}
\end{equation}
The two tie together through the loop of \cref{fig:loop}:
\begin{enumerate}[nosep]
\item Collision pulls populations toward the local equilibrium ($f^{\mathrm{eq}}$ shares
      $\rho$ and $\rho\vu$ with $f$, so mass and momentum are conserved locally).
\item Streaming moves those relaxed populations to neighbouring nodes, so each node now
      holds a new mix and therefore new moments.
\item Collision recomputes $f^{\mathrm{eq}}$ from those new moments. $f^{\mathrm{eq}}$ is
      never stored across steps; it is regenerated from the streamed state every time.
\end{enumerate}

\begin{figure}[h]
\centering
\begin{adjustbox}{max width=\linewidth}% The collide-stream loop and where the collision seam sits.
\begin{tikzpicture}[
  >={Stealth[length=2.2mm]},
  box/.style={draw, rounded corners=3pt, align=left, inner sep=7pt, font=\small},
  col/.style={box, fill=orange!8, draw=orange!70!black, text width=100mm},
  str/.style={box, fill=blue!7, draw=blue!60!black, text width=100mm},
  bnd/.style={box, fill=gray!8, draw=gray!60, text width=100mm},
  arr/.style={->, line width=1.1pt},
  note/.style={font=\small, gray!80!black, align=left},
  seam/.style={densely dashed, line width=1.1pt, red!70!black},
]
% collision
\node[col] (col) at (0,0) {%
  \textbf{Collision} \hfill local, per node\\[3pt]
  \begin{tabular}{@{}l@{\qquad}l@{}}
  {\scriptsize\ttfamily rho, u = moments(lattice, f)} & {\footnotesize $\rho=\sum_i f_i,\quad \rho\vu=\sum_i \vc_i f_i$}\\[2pt]
  {\scriptsize\ttfamily feq = equilibrium(lattice, rho, u)} & {\footnotesize $f^{\mathrm{eq}}(\rho,\vu)$}\\[2pt]
  {\scriptsize\ttfamily f\_star = f - (f - feq) / tau} & {\footnotesize $f_i^{*} = f_i - \tfrac{1}{\tau}(f_i - f_i^{\mathrm{eq}})$}
  \end{tabular}};
% streaming
\node[str, below=24mm of col] (str) {%
  \textbf{Streaming} \hfill non-local, a permutation\\[3pt]
  \begin{tabular}{@{}l@{\qquad}l@{}}
  {\scriptsize\ttfamily f = stream(lattice, f\_star)} & {\footnotesize $f_i(\mathbf{x}+\vc_i\Delta t,\,t+\Delta t)=f_i^{*}(\mathbf{x},t)$}
  \end{tabular}};
% boundaries
\node[bnd, below=8mm of str] (bnd) {%
  \textbf{Boundaries} \hfill overwrite the wrapped entries\\[3pt]
  \begin{tabular}{@{}l@{\qquad}l@{}}
  {\scriptsize\ttfamily f = bc(f)} & {\footnotesize bounce-back, Zou--He, \ldots}
  \end{tabular}};

% loop arrows
\draw[arr] (col.south) -- (str.north);
\node[note, anchor=north west, text width=64mm] at ($(col.south)!0.45!(str.north)+(0.3,0)$)
  {$f^{*}$: relaxed populations; $f - f^{\mathrm{eq}}$ rides along to the neighbours};
\draw[arr] (str.south) -- (bnd.north);
\draw[arr, rounded corners=4pt] (bnd.west) -- ++(-1.3,0) |- (col.west);
\node[note, anchor=east, text width=36mm, align=right] at ($(col.west)!0.5!(bnd.west)+(-1.45,0)$)
  {new mix at every node, so new moments; $f^{\mathrm{eq}}$ is rebuilt, never stored};

% the seam
\coordinate (sl) at ($(col.south west)!0.3!(str.north west)+(-2.0,0)$);
\coordinate (sr) at ($(col.south east)!0.3!(str.north east)+(6.2,0)$);
\draw[seam] (sl) -- (sr);
\node[note, red!70!black, anchor=south west, text width=56mm] at ($(col.south east)+(0.35,0.1)$)
  {\textbf{the seam} (\code{CollisionOperator}): everything that depends on $f^{\mathrm{eq}}$ sits above; a learned or quantum operator replaces that box};
\node[note, blue!60!black, anchor=north west, text width=56mm] at ($(str.north east)+(0.35,-0.05)$)
  {fixed, trivially differentiable permutation; a permutation unitary on the quantum track};
\node[note, anchor=north west, text width=56mm] at ($(bnd.north east)+(0.35,-0.05)$)
  {$\nu = \cs^2(\tau - \tfrac12)\Delta t$: the $\tfrac12$ is the one-step discrete shift};
\end{tikzpicture}
\end{adjustbox}
\caption{The collide--stream loop and the seam. Collision is local and is the only place
$f^{\mathrm{eq}}$ is built; streaming is a fixed permutation that carries the relaxed
populations, and with them $f - f^{\mathrm{eq}}$, to the neighbours; boundaries then
overwrite the wrapped entries they own. Everything a learned or quantum operator would
replace sits above the dashed line.}
\label{fig:loop}
\end{figure}

The physics of the coupling: $f^{\mathrm{eq}}$ carries the inviscid transport (its second
moment is the Euler pressure tensor, which is why the truncation order of $f^{\mathrm{eq}}$
in $\vu$ matters, \cref{sec:eq-theory}). The non-equilibrium part $f - f^{\mathrm{eq}}$,
which streaming carries to the neighbours before the next relaxation, is what produces
viscous stress. That is where
\begin{equation}
\nu = \cs^2\left(\tau - \tfrac{1}{2}\right)\Delta t
\label{eq:viscosity}
\end{equation}
comes from \parencite{Kruger2017}, and the $\tfrac{1}{2}$ is purely a consequence of
streaming being a discrete one-step shift.

That is also why the seam in \code{fluxlb} is drawn where it is: streaming is a fixed,
trivially differentiable permutation; everything that depends on $f^{\mathrm{eq}}$, and
hence everything a learned or quantum operator would replace, sits inside the collision
module.

%----------------------------------------------------------------------------------------

\section{Implementation}\label{sec:stream-impl}

\paragraph{Convention.}
Push, functional, periodic. \code{stream(lattice, f)} returns a new tensor in which slab $q$
has been rolled by $+\code{lattice.shifts[q]}$, so after the call every node holds the
populations that arrived there. The function has no state and knows nothing about walls: a
non-periodic boundary is applied after it and overwrites the wrapped values it owns.

\paragraph{\code{stream(lattice, f)}.}
Takes $f$ of shape $[Q, \ast]$ with $\code{f.ndim} = D + 1$ and returns the streamed field
with the same shape, dtype and device. \code{ValueError} is raised if \code{f.shape[0]}
differs from \code{lattice.Q} or \code{f.ndim} from \code{lattice.D + 1}; the checks run
before any allocation. The body is one roll per direction and one stack
(\cref{fig:tensor}):
\begin{lstlisting}[style=py]
dims = tuple(range(lattice.D))
return torch.stack(
    [torch.roll(f[q], shifts=lattice.shifts[q], dims=dims) for q in range(lattice.Q)]
)
\end{lstlisting}

\begin{figure}[h]
\centering
\begin{adjustbox}{max width=\linewidth}% The tensor view of stream(): per-direction roll of a [*spatial] slab, then stack.
\begin{tikzpicture}[
  >={Stealth[length=2.2mm]},
  slab/.style={draw, fill=blue!7, draw=blue!60!black, minimum width=22mm, minimum height=15mm, font=\small\ttfamily},
  oslab/.style={slab, fill=green!8, draw=green!50!black},
  arr/.style={->, line width=1pt},
  note/.style={font=\small, gray!80!black, align=center},
  lbl/.style={font=\small, inner sep=1pt},
]
% input stack f[Q, Nx, Ny]: back slabs up and to the right, front slab drawn last
\foreach \k in {3,2,1,0} {
  \node[slab] (in\k) at ({0.25*\k}, {0.25*\k}) {};
}
\node[lbl] at (in0) {f[q]};
\node[lbl, anchor=south] at ($(in3.north)+(-0.3,0.12)$) {$f\,[Q, N_x, N_y]$};
\node[note, anchor=north, text width=42mm] at ($(in0.south)+(0.4,-0.25)$)
  {one slab per direction;\\ \code{f[q]} has axes $0,\dots,D-1$};

% roll
\node[draw, rounded corners=2pt, fill=white, font=\small\ttfamily, align=left, anchor=west, inner sep=6pt] (roll) at (3.4,0)
  {torch.roll(f[q],\\ \ shifts=lattice.shifts[q],\\ \ dims=range(D))};
\draw[arr] (in0.east) -- (roll.west);
\node[note, anchor=north, text width=62mm] at ($(roll.south)+(0,-0.3)$)
  {\code{shifts[q]} $= \vc_q$ as host integers:\\ no device sync, every axis wraps};

% output stack
\foreach \k in {3,2,1,0} {
  \node[oslab] (out\k) at ({13.0+0.25*\k}, {0.25*\k}) {};
}
\node[lbl] at (out0) {roll(f[q])};
\node[lbl, anchor=south] at ($(out3.north)+(-0.3,0.12)$) {$f'\,[Q, N_x, N_y]$};
\draw[arr] (roll.east) -- node[above=3pt, lbl, font=\small\ttfamily] {torch.stack} (out0.west);
\node[note, anchor=north, text width=46mm] at ($(out0.south)+(0.4,-0.25)$)
  {a new tensor; $f$ is untouched,\\ so no in-place write enters autograd};
\end{tikzpicture}
\end{adjustbox}
\caption{The tensor view. Each slab \code{f[q]} has already dropped the direction axis, so
its spatial axes are $0,\dots,D-1$ and the roll pairs axis $i$ with component $i$ of
$\vc_q$. The rolls allocate fresh tensors and \code{torch.stack} assembles them; nothing is
written into a tensor that autograd holds.}
\label{fig:tensor}
\end{figure}

\paragraph{Axis pairing.}
The population tensor pairs spatial axis $i$ with \code{dim=i + 1}
(\cref{Chapter1}), but a slab \code{f[q]} pairs it with \code{dim=i}. Passing the
full-tensor dims to the per-slab roll is an off-by-one that fails on every lattice with an
\code{IndexError}, which is why the implementation binds \code{dims} once from
\code{lattice.D} rather than from \code{f.ndim}.

\paragraph{Why stack, not assign.}
The alternative of cloning \code{f} and assigning \code{f[q] = torch.roll(...)} slab by slab
is an in-place write on a tensor in the autograd graph. It passes a gradient check as long
as nothing has saved a view of the clone, and fails with \emph{``a variable needed for
gradient computation has been modified by an inplace operation''} as soon as something has.
The stack form has no such dependence on the caller and is one copy cheaper: $Q$ rolls plus
one stack, against a clone, $Q$ rolls and $Q$ slab copies.

\paragraph{Host-side shifts.}
\code{lattice.shifts} is a plain Python list of integer tuples built once in
\code{Lattice.\_\_init\_\_} (\cref{sec:lattice-impl}), so \code{torch.roll} receives its
arguments without reading a device tensor and the step never synchronises with the host.

\paragraph{Memory.}
Per call, $Q$ slab-sized temporaries and one $[Q, \ast]$ output are allocated; the input is
untouched. In a differentiable rollout the output is what autograd retains, since the
backward of a permutation needs no saved input, so the per-step cost of streaming is one
population field. Checkpointing across steps (roadmap E0.2) bounds the total.

\paragraph{Differentiability.}
The Jacobian of \cref{eq:stream} is a permutation matrix, so the adjoint is the inverse
permutation: the gradient of a loss with respect to $f$ is the upstream gradient streamed
with every shift negated. The tests check this identity exactly (\cref{Chapter6}).

\paragraph{Usage.}
\begin{lstlisting}[style=py]
from fluxlb.core.lattices import D2Q9
from fluxlb.core.streaming import stream

lat = D2Q9()
f = stream(lat, f_star)        # push, periodic wrap; f_star untouched
stream(lat, f_star[:5])        # ValueError: f.shape[0] (5) != lattice.Q (9)
\end{lstlisting}
